% Source: https://artofproblemsolving.com/community/c3h2071810_infinite_roots
% Stable identity: aops-topic-2071810-q01
Let \(y=\sqrt{2+\sqrt{2+\sqrt{2+\cdots}}}\), the limit of \(y_1=\sqrt2\), \(y_{n+1}=\sqrt{2+y_n}\). If \(y_n<2\) then \(y_{n+1}<\sqrt4=2\), and \(y_{n+1}>y_n\iff 2+y_n>y_n^2\iff (2-y_n)(1+y_n)>0\), which holds. So \((y_n)\) is increasing and bounded by \(2\), hence converges to some \(y\) with \(y=\sqrt{2+y}\). Then \(y^2-y-2=(y-2)(y+1)=0\), and since \(y>0\), \(y=2\).

Therefore \(x=\sqrt{1+y}=\boxed{\sqrt3}\). (The poster's approach of squaring the whole expression stalls because the \(1\) breaks the self-similarity; the repeating part is the inner radical \(y\).)
